\documentclass{beamer}

\mode<presentation>
{
  \usetheme{Singapore}      % or try Darmstadt, Madrid, Warsaw, ...
  \usecolortheme{crane} % or try albatross, beaver, crane, ...
  \usefonttheme{serif}  % or try serif, structurebold, ...
  \setbeamertemplate{navigation symbols}{}
  \setbeamertemplate{caption}[numbered]
} 

\usepackage[autostyle=true]{csquotes}
\usepackage{xeCJK}
\usepackage[T1]{fontenc}
\usepackage{xcolor}
\usepackage{setspace}
\setmainfont[AutoFakeSlant=0.2]{Charis SIL}
\setCJKmainfont[AutoFakeSlant=0.2]{Noto Serif CJK HK}
\setCJKmonofont[AutoFakeSlant=0.2]{Noto Serif CJK HK}
\usepackage{polyglossia}
\newfontfamily\thaifont[Script=Thai]{Noto Serif Thai}
\setmainlanguage{english}
\setotherlanguages{thai}
\usepackage{csquotes}
\usepackage{array}
\usepackage{makecell}
\usepackage{changepage}
\usepackage{textcomp}

\title[Pearl River Sinitic]{Sinitic Languages of the Pearl River Region \\ --- Unit 3.3}
\subtitle{13ᵗʰ EACL Summer School, Universität Stuttgart}
\author{Hilário de Sousa 蘇沙}
\institute{EHESS–CRLAO \\ {\tiny membre associé}}
\date{Wednesday 29 July 2026 \\ 丙午六月十六 \\ {\footnotesize \texttt{hilario@bambooradical.com}}}

\begin{document}

\begin{frame}
  \titlepage
\end{frame}

\section{Kra-Dai influence}

\begin{frame}[fragile]{Non-Sinitic influence}

{\small Ethnolinguistic situation before the Hàn expansion was something like this:\\}

    \centering
    \includegraphics[width=0.6\linewidth]{southchina}

{\flushleft
\vfill \setstretch{0.5} {\tiny 
  Li, Hui et al. 2007. Mitochondrial DNA diversity and population differentiation in southern East Asia. \textit{American Journal of Physical Anthropology} 134(4). 481–488. \url{https://doi.org/10.1002/ajpa.20690} \url{https://loca.fudan.edu.cn/Doc/A35.pdf}\par}}

\end{frame}

\begin{frame}[fragile]{Non-Sinitic influence}

Sino-Tibetan 漢藏 / Trans-Himalayan 跨喜瑪拉雅

\begin{columns}
\begin{column}{0.9\textwidth}
   \includegraphics[width=1\linewidth]{sinotibetan}
\end{column}
\begin{column}{0.15\textwidth}
    {\scriptsize All non-Sinitic Sino-Tibetan languages: \enquote{Tibeto-Burman} (one node? multiple nodes?) Is Sinitic a first order node?}\\
\end{column}
\end{columns}

\end{frame}

\begin{frame}[fragile]{Non-Sinitic influence}

\begin{columns}[t]
\begin{column}[t]{0.3\textwidth}
Austroasiatic 南亞\\
\end{column}
\begin{column}[t]{0.7\textwidth}
    \flushright
    \includegraphics[width=1\linewidth]{austroasiatic}
\end{column}
\end{columns}

\end{frame}

\begin{frame}[fragile]{Non-Sinitic influence}

\begin{columns}[t]
\begin{column}[t]{0.3\textwidth}
Hmong-Mien 苗瑤\\
\end{column}
\begin{column}[t]{0.7\textwidth}
    \flushright
    \includegraphics[width=1\linewidth]{hmongmien}
\end{column}
\end{columns}

\end{frame}

\begin{frame}[fragile]{Non-Sinitic influence}

Austronesian 南島
\bigskip

    \includegraphics[width=1\linewidth]{austronesian}

\end{frame}

\begin{frame}[fragile]{Kra-Dai influence}

\begin{columns}[t]
\begin{column}[t]{0.5\textwidth}
Kra-Dai {\scriptsize (Tai-Kadai)} 壯侗／侗台\\
{\scriptsize
\begin{itemize}
    \item Tai 台／壯傣
    \item Be {\scriptsize Ong-Be} 臨高
    \item Kam-Sui 侗水
    \item Hlai 黎
    \item Kra 仡央 (仡佬布央)
    \item Biao–Lakkja 標–拉珈
\end{itemize}}

\flushleft
    \includegraphics[width=1\linewidth]{kradailiao2023}
{\vfill \setstretch{0.5} {\tiny Liao, Hanbo. 2023. Tonal Behavior as of Areal and Typological Concerns: Centering on the Sinitic and Kam-Tai Languages in Lingnan. \textit{Languages} 8(2). 148. \url{https://doi.org/10.3390/languages8020148}}\par}

\end{column}
\begin{column}[t]{0.5\textwidth}
    \flushright
    \includegraphics[width=1\linewidth]{kradaimap}
\end{column}
\end{columns}

\end{frame}

\begin{frame}[fragile]{Kra-Dai influence}

\begin{columns}
    \begin{column}{0.4\linewidth}
    那\_\_ placenames\\
    {\scriptsize (Not Mandarin 那 nà \textsuperscript{tone C} or Cantonese naa5 \textsuperscript{tone B2})}\\
    \medskip
    那\_\_ placenames:\\ naa4 \textsuperscript{tone A2}
    {\tiny \setstretch{0.5} \url{https://www.youtube.com/watch?v=A8GCPYdRmzM}}\\
    \medskip
     {\scriptsize or no4 in 四邑,\\ e.g. 那琴 Canto. no4 kam4\\
    {\tiny \setstretch{0.5} \url{https://youtu.be/C6foBnDLN48?si=dcdGJ6RIVSRM357M&t=12}}}\\
    Proto-Tai *na\textsuperscript{A} \enquote*{paddy field}\\
    e.g. 西雙版納 [ɕip˧˥ sɔŋ˥ pan˥˧ naː˥˧] \enquote*{12k fields}
    \\
    \end{column}
    \begin{column}{0.6\linewidth}
    \includegraphics[width=1\linewidth]{naplacenames}
    \medskip
   
    \end{column}
\end{columns}

\end{frame}

\begin{frame}[fragile]{Kra-Dai influence}

{\small Many Kra-Dai words in Standard Cantonese basic vocabulary e.g.:}\\
{\tiny Tai examples given out of convenience; donor language can be from another branch of Kra-Dai}\\
\medskip
{\footnotesize
\begin{tabular}{l|l}
     Stan. Canto. & Kra-Dai forms\\ \hline
     han4 \textsuperscript{A2} \enquote*{to itch} & PT *ɣal\textsuperscript{A} \enquote*{to itch}, 北壯 han\textsuperscript{A2}, 泰 \begin{thai}คัน\end{thai} kʰan\textsuperscript{A2}\\
     naat3 \textsuperscript{DL1} \enquote*{to scorch} & PT *ɗɤːt\textsuperscript{D} \enquote*{to boil \textsuperscript{intr}}, 布依 daːt\textsuperscript{DL1}, 泰 \begin{thai}เดือด\end{thai} dɯət\textsuperscript{DL1}\\
     luk6 \textsuperscript{D2} \enquote*{to scald} & PT luək\textsuperscript{D} \enquote*{to scald}, 龍州壯 luːk\textsuperscript{DL2}, 泰 \begin{thai}ลวก\end{thai} luək\textsuperscript{DL2}\\
     kam2 \textsuperscript{B1} \enquote*{to cover} & PB *kʰum\textsuperscript{BC1}; PT *hɤm\textsuperscript{B} \enquote*{cover w/ cloth}, 泰 \begin{thai}ห่ม\end{thai} hom\textsuperscript{B1}\\
     lam3 \textsuperscript{C1} \enquote*{to collapse} & PT *lɤm\textsuperscript{C} \enquote*{to fall down}, 布依 lam\textsuperscript{C2}, 泰 \begin{thai}ล้ม\end{thai} lom\textsuperscript{C2}\\
     tap6 \textsuperscript{D2} \enquote*{to pound} & PT *dup\textsuperscript{D} \enquote*{to pound}, 布依 tup\textsuperscript{DS2}, 泰 \begin{thai}ทุบ\end{thai} tʰup\textsuperscript{DS2}\\
     dan2 \textsuperscript{B1} \enquote*{to put (down)} & Pre-Hlai *ʈəːnɦ \enquote*{to put (down)}\\
     ngaam1 \textsuperscript{A1} \enquote*{correct} & 北壯 ŋaːm\textsuperscript{B1} \enquote*{exact}, PSWT ŋaːm\textsuperscript{A2} \enquote*{nice}, 泰 \begin{thai}งาม\end{thai} ŋaːm\textsuperscript{A2}\\ 
     hai1 \textsuperscript{A1} \enquote*{vagina} & PT 
     *hiː\textsuperscript{A} \enquote*{vagina}, 上思壯 hoɥ\textsuperscript{A1}, 泰 \begin{thai}หี\end{thai} hiː\textsuperscript{A1}\\
     ke1 \textsuperscript{A1} \enquote*{excrement} & PT * C̬.qɯj\textsuperscript{C}, 上思壯 kʰoɥ\textsuperscript{C1}, 泰 \begin{thai}ขี้\end{thai} kʰiː\textsuperscript{C1}\\
      & PB *kaːj\textsuperscript{BC2}; PKS *ke\textsuperscript{4}; PKr *kai\textsuperscript{C}; PHl *ɦaːjʔ\\
     %gaap3 \textsuperscript{DL1} gwaai2 \textsuperscript{B1} \enquote*{frog} & PT *kɤp\textsuperscript{D} \enquote*{frog}, 布依 kap\textsuperscript{DS1}, 泰 \begin{thai}กบ\end{thai} kop\textsuperscript{DS1}; PB *kup\textsuperscript{D2} \enquote*{big frog}\\
     % & PT *krweː\textsuperscript{C}, 北壯 kweː\textsuperscript{C1}\\
\end{tabular}}
\medskip

{\footnotesize Notice: Normally, nasal/liquid onset with high tone cannot be a reflex of Middle Chinese words}

\end{frame}

\begin{frame}[fragile]{Kra-Dai influence}

{\small Tons of Zhuàng words in Nánníng Pínghuà \& Nán. Canto. e.g.:}\\
{\footnotesize
\begin{tabular}{l|l|l}
    南寧位子淥平 & 南寧白 & 壯(邕) \\ \hline
    hɐj˧˥ \textsuperscript{C1a} \enquote*{to give} & & 北 haɰ˥ \textsuperscript{C1}\\ 
    (心墟平 həɰ˧˥ \textsuperscript{C1a}) & & 南 haj˥˧ \textsuperscript{B2}\\
    (亭子平 həj˥ \textsuperscript{C1}) & & (PT *haɰ\textsuperscript{C} \enquote*{to give})\\ \hline
    ɲaːp˥ \textsuperscript{D*} \enquote*{vexed} & & ɲaːp\textsuperscript{DL1} \enquote*{irritated}\\
     & & (PT *ʰɲaːp\textsuperscript{D} \enquote*{coarse})\\ \hline
    ɬɔːt˧ \textsuperscript{D1} \enquote*{buttocks} & & 北 ɣot\textsuperscript{}{DS1}\\
     & & 南 tot\textsuperscript{DS1}\\
     & & (PT *k.tɤt\textsuperscript{D} \enquote*{fart})\\ \hline 
    %nɐm˥˧ \textsuperscript{A1} \enquote*{stupid} & nɐm˥ \textsuperscript{A1} \enquote*{unsophisticated} & 北 nam\textsuperscript{A1} \enquote*{slow thinking}\\ \hline
    -tɪk˥ \textsuperscript{D*} \enquote*{place} & -tiː˥ \textsuperscript{A1} \enquote*{place} & 北 tiːᵊk\textsuperscript{D2} / tiː\textsuperscript{B2}\\
     & & (PT *diː\textsuperscript{B} \enquote*{place})\\ \hline
    nɛːp˥ \textsuperscript{D*} \enquote*{to hold btw} & nɛːp˧ \textsuperscript{DL1} \enquote*{to hold btw} & 北 neːp\textsuperscript{DL1}\\
     & & (PT *ʰniːp\textsuperscript{D} \enquote*{to pinch})\\ \hline
    nɛːw˥˧ \textsuperscript{A1} \enquote*{sticky} & niːw˥ \~ nɛːw˥ \textsuperscript{A1} & 北 niːᵊw\textsuperscript{A1}\\
     & & 南 niːw\textsuperscript{A1us}\\
     & & (PT *ʰniəw\textsuperscript{A} \enquote*{sticky})\\ %\hline
    %kʷɛː˨˩ \textsuperscript{A2}(-tʃiː˧ 子) \enquote*{小 frog} & kʷʰɛː˧˥ \textsuperscript{B1} \enquote*{小 frog} & 北 kweː\textsuperscript{C1}\\
    % & & 南 kʰleː\textsuperscript{C1}\\
    % & & (PT *krweː\textsuperscript{C} \enquote*{little frog})\\
\end{tabular}}

\end{frame}


\end{document}
